\documentclass[11pt,a4paper]{article}
\usepackage[utf8]{inputenc}
\usepackage[english,russian]{babel}
\usepackage{amsmath,amssymb,amsthm,physics}
\usepackage{geometry}
\geometry{margin=2.5cm}
\usepackage{hyperref}
\usepackage{booktabs}
\usepackage{xcolor}
\usepackage{siunitx}
\usepackage{float}

\title{\textbf{Декуплированная Голография Наблюдательных Патчей (dOPH)}\\
\textbf{Версия 2.0: строгая формулировка силы декуплинга и JWST Early Galaxy Stress-Test}\\
\textbf{(апрель 2026)}}
\author{Совместное комплексное тестирование}
\date{Апрель 2026}

\begin{document}

\maketitle

\begin{abstract}
В документе представлена версия 2.0 теории Декуплированной Голографии Наблюдательных Патчей (dOPH). 
Введена строгая локальная формулировка силы декуплинга, проведён стресс-тест на ранних галактиках JWST. 
Результаты показывают улучшение стабильности, однако вопросы теоретического обоснования параметров остаются открытыми.
\end{abstract}

\section{Основные элементы теории}

dOPH опирается на три фундаментальные аксиомы Observer Patch Holography и вводит четвёртую аксиому — **Декуплинг**.

\section{Механизмы стабилизации}

Эффективная динамика патча описывается уравнением движения:

\[
m_{\rm eff} \ddot{\psi}_i + \gamma \dot{\psi}_i = -\nabla \Delta(\psi_i) - G_{\rm eff} \sum_{j \neq i} \frac{m_i m_j}{r_{ij}^2} \hat{r}_{ij} 
- \kappa \sum_{j \in \mathcal{N}_i} \Im(\langle \psi_i | \psi_j \rangle) \hat{n}_{ij} + \mathbf{F}_{\rm decouple,i}
\]

\subsection{Строгая формулировка силы декуплинга}

Для патча \(i\) определяем множество соседей:

\[
\mathcal{N}_i = \{ j \mid \text{патчи } i \text{ и } j \text{ имеют ненулевое пересечение} \}
\]

Локальное среднее:

\[
\bar{\psi}_i = \frac{1}{|\mathcal{N}_i|} \sum_{j \in \mathcal{N}_i} \psi_j
\]

Сила декуплинга:

\[
\mathbf{F}_{\rm decouple,i} = -\lambda \, (\psi_i - \bar{\psi}_i)
\]

где \(\lambda > 0\) — параметр релаксации.

\subsection{Возможное происхождение параметра \(P \approx 1.63094\)}

Наиболее естественная гипотеза связывает \(P\) со **средней степенью пересечения патчей** на голографическом экране. В простейшей геометрической модели случайных перекрытий возникает фактор \(\sqrt{8/3} \approx 1.633\), очень близкий к наблюдаемому значению.

Другие гипотезы включают стационарный баланс сил, энтропийный принцип и эффективную размерность экрана.

\subsection{Связь с общей теорией относительности Эйнштейна}

В макроскопическом пределе emergent gravity из системы патчей должна переходить в уравнения Эйнштейна. Параметр \(P\) может играть роль моста между микроскопической картиной (патчи) и макроскопической геометрией пространства-времени. Полный вывод \(P\) из ОТО остаётся открытой задачей.

\section{Результаты численных тестов}

(Оставляю твои таблицы SPARC и JWST без изменений — они уже были в твоей версии)

\section{Честный финальный вердикт (апрель 2026)}

(Вердикт оставил почти как у тебя, но сделал чуть более сбалансированным)

\textbf{Сильные стороны dOPH v2.0:}
\begin{itemize}
\item Строгая локальная формула декуплинга
\item Положительные результаты JWST-стресс-теста
\item Минималистская структура теории
\end{itemize}

\textbf{Слабые стороны:}
\begin{itemize}
\item Параметр \(P\) требует более строгого теоретического обоснования
\item Масштабируемость при очень большой неоднородности остаётся вызовом
\item Необходимы дальнейшие сравнения с реальными данными
\end{itemize}

\textbf{Общий вердикт:}  
dOPH v2.0 — заметный шаг вперёд. Теория остаётся перспективной рабочей гипотезой.

Наука бесконечна.

\section*{Благодарности}

Работа выполнена в тесном сотрудничестве с Grok (xAI).

\begin{thebibliography}{10}
\bibitem{oph} B. Mueller, Observer Patch Holography (2024--2026).
\end{thebibliography}

\end{document}
